\chapter{Literature Review and Related Work}
\label{chap:relatedworks}

\section{Competitor Analysis}
\label{section:competitor-analysis}

\begin{figure}[h]
    \centering
    \includegraphics[width=0.8\textwidth]{chapter-2/competitor-landscape-analysis.png}
    \caption{Competitive landscape positioning based on software granularity specialization and
    multi-language support.}
\end{figure}

To evaluate the current state of software code smell detection tools, a competitive analysis was conducted using the
Asana Competitor Analysis Framework. The analysis focused on tools specializing at different code granularity levels and
their multi-language capabilities. We compare each solution by the following features:
\textbf{1.} Analysis granularity level
\textbf{2.} Analysis capabilities
\textbf{3.} Supported languages
\textbf{4.} Report format

\begin{table}[htbp]
    \centering
    \begin{tabularx}{\textwidth}{| p{3.5cm} | X | X | X | X |}
        \hline
        \textbf{Features} & \textbf{Designnite} & \textbf{Sonarqube} & \textbf{Go architect} & \textbf{ArchDiver} \\
        \hline
        Level of Interest 
        & Code + Architecture Level 
        & Code-level Design 
        & Architecture Level
        & Architecture Level \\
        \hline
        Analysis 
        & Architecture Smells, Architecture Metrics 
        & Security / Vulnerability smells, Code Smells
        & Architecture Metrics  
        & Architectural Smell, Architecture Metrics \\
        \hline
        Language 
        & Java, C\#, Python 
        & Supports 35 Programming Languages
        & Golang 
        & Language Independent Analysis \\
        \hline
        Explanation 
        & Architectural Metrics Report, Detected Smells
        & Security Metrics, Code quality Metrics
        & Architectural Metrics Report
        & Architectural Metrics, Detected Smells, Natural Language Explanation\\
        \hline
    \end{tabularx}
    \caption{Comparison of software analysis tools by feature set}
    \label{tab:tool_comparison}
\end{table}

From the comparison table, we can see that exsisting competitors...

\section{Literature Review}
\label{section:literature-review}

Tushar et al.[1] define architecture smells as "indicators of poor quality"
manifesting within source code at an architectural level of granularity.
He explores the association between architecture smells and design smells
in order to find a correlation between smells in different levels of granularity.
He concludes that while architecture and design smells have a general correlation,
no significant link is established between specific design and architecture smells.
Although the paper explores the relation between architecture and design smells,
it does not explore the classification of architecture smells.
\\

Hanyu et al.[2] devised an architecture for detecting code smells using 3 GNN models, GCN, GAT, and
GraphSAGE for detecting and refactoring large class, Long method, and Feature envy code smells. 
Using a combination of a method level Program Dependency Graph (PDG) and a class level Direct Heterogenous Graph (DHG).
They construct the graph using ... . The proposed model has an average accuracy of 42 percent
with GraphSAGE having the best accuracy out of the 3 tested models. They summarize that ... .
In this work, we aim to expand on their literature by focusing on architecture smells.
\\

Smajevic et al.[3] proposed a method for detecting Enterprise architecture smells using Knowledge graphs and GNN models.
He .... 